\documentclass[aspectratio=169]{beamer}
\usetheme{Madrid}
\usecolortheme{default}
\setbeamertemplate{navigation symbols}{}
\usepackage[T1]{fontenc}
\usepackage[utf8]{inputenc}
\usepackage{lmodern}
\usepackage{amsmath,amssymb}
\usepackage{booktabs}
\usepackage{array}
\usepackage{listings}
\usepackage{xcolor}
\usepackage{hyperref}

\definecolor{codebg}{RGB}{247,247,247}
\definecolor{darkblue}{RGB}{20,50,120}
\definecolor{darkgreen}{RGB}{0,110,60}
\definecolor{darkred}{RGB}{140,30,30}

\lstset{
  basicstyle=\ttfamily\scriptsize,
  backgroundcolor=\color{codebg},
  frame=single,
  framerule=0.2pt,
  rulecolor=\color{black!30},
  breaklines=true,
  columns=fullflexible,
  keywordstyle=\color{darkblue},
  commentstyle=\color{darkgreen},
  stringstyle=\color{darkred},
  showstringspaces=false
}

\title{Using the Python ILP Matching Program}
\author{Quick guide for tuning and interpreting results}
\date{}

\begin{document}

\begin{frame}
  \titlepage
\end{frame}

\begin{frame}{What the program does}
\begin{itemize}
  \item Reads one Excel sheet with mixed respondents: \textbf{students} and \textbf{locals}.
  \item Builds feasible student--local pairs using \textbf{hard constraints}.
  \item Computes a weighted compatibility score for each feasible pair.
  \item Solves a \textbf{one-to-one Integer Linear Program}.
  \item Default objective is \textbf{lexicographic}:
  \begin{enumerate}
    \item maximize number of matches,
    \item then maximize the \textbf{minimum selected score} (Chebychev idea),
    \item then slightly prefer larger total score.
  \end{enumerate}
\end{itemize}
\vspace{0.3em}
\textbf{Main outputs:} \texttt{matches.csv} and \texttt{*.diagnostics.csv}
\end{frame}

\begin{frame}{Optimization model in one slide}
Binary variable:
\[
 x_{ij}=\begin{cases}
 1 & \text{if student } i \text{ is matched with local } j,\\
 0 & \text{otherwise.}
 \end{cases}
\]
Chebychev helper variable:
\[
 t \ge 0 \quad \text{(minimum score among selected pairs)}
\]
Constraints:
\[
\sum_j x_{ij} \le 1 \quad \forall i, \qquad \sum_i x_{ij} \le 1 \quad \forall j
\]
\[
 t \le s_{ij} + (1-x_{ij}) \quad \forall (i,j) \text{ feasible}
\]
Default solve order:
\begin{enumerate}
  \item maximize $\sum_{i,j} x_{ij}$,
  \item among those, maximize $t$,
  \item add a tiny tie-break on $\sum s_{ij}x_{ij}$.
\end{enumerate}
\end{frame}

\begin{frame}{Hard constraints: where to change them}
The hard constraints are handled before the ILP by three functions:
\begin{itemize}
  \item \texttt{gender\_compatible(...)}
  \item \texttt{pets\_compatible(...)}
  \item \texttt{must\_match\_ok(...)}
\end{itemize}
A pair is feasible only if \textbf{all} of these return \texttt{True} in the required directions.

\vspace{0.4em}
\textbf{Typical edits:}
\begin{itemize}
  \item Make a rule stricter: return \texttt{False} for more cases.
  \item Make a rule softer: treat more patterns as \texttt{no preference}.
  \item Temporary debugging: use CLI flags
  \texttt{--relax-gender}, \texttt{--relax-pets}, \texttt{--relax-mustmatch}.
\end{itemize}
\end{frame}

\begin{frame}[fragile]{Example: changing hard-constraint boundaries}
\textbf{Gender / pets / must-match are rule-based.}
Change the boundary by editing the condition logic, for example:

\begin{lstlisting}[language=Python]
def pets_compatible(a_pets, b_pets, relax=False):
    if relax:
        return True
    if has_no_pref(a_pets) or has_no_pref(b_pets):
        return True

    bad_words = {"severe allergies", "cannot have pets"}
    has_pets_words = {"has cat", "has dog", "has pets"}

    if (norm_set(a_pets) & bad_words and norm_set(b_pets) & has_pets_words):
        return False
    return True
\end{lstlisting}

To soften the boundary, reduce \texttt{bad\_words}. To tighten it, add more conflict keywords.
\end{frame}

\begin{frame}{Chebychev weights: where to change them}
The weighted score is computed from named components:
\begin{center}
\begin{tabular}{lr}
\toprule
Component & Default weight \\
\midrule
Profile & 34 \\
Hobbies & 23 \\
Appreciate & 10 \\
Languages & 9 \\
Age & 9 \\
Info & 7 \\
Degree & 5 \\
Field & 3 \\
\bottomrule
\end{tabular}
\end{center}

In code, edit:
\begin{itemize}
  \item \texttt{DEFAULT\_WEIGHTS = \{...\}}
  \item or override at runtime with \texttt{--weights}
\end{itemize}
The program normalizes the numbers internally, so they only need to be \textbf{relative priorities}.
\end{frame}

\begin{frame}[fragile]{Example: changing weights}
\textbf{Inside the script:}
\begin{lstlisting}[language=Python]
DEFAULT_WEIGHTS = {
    "profile": 34,
    "hobbies": 23,
    "appreciate": 10,
    "languages": 9,
    "age": 9,
    "info": 7,
    "degree": 5,
    "field": 3,
}
\end{lstlisting}

\textbf{At the command line:}
\begin{lstlisting}[language=bash]
python main.py --excel "Refined Attributes with options.xlsx" \
  --weights=hobbies=30,profile=28
\end{lstlisting}

This keeps the same model but shifts relative emphasis before normalization.
\end{frame}

\begin{frame}{How the input matrix is specified}
The program expects one Excel sheet, usually named:
\[
\texttt{the refined data}
\]
Important columns include:
\begin{itemize}
  \item IDs and status: \texttt{Identification}, \texttt{My status}
  \item matching attributes: profile, hobbies, languages, age, pets
  \item study descriptors: degree and field
  \item must-match fields: criteria 1/2 and their specs
\end{itemize}

\textbf{Important preprocessing:}
\begin{itemize}
  \item remove rows marked \texttt{REMOVED}
  \item drop the legend/options row below the header
  \item normalize spaces and multi-value separators
\end{itemize}
\end{frame}

\begin{frame}[fragile]{Minimal input structure}
Each row is one person. The solver first classifies rows as \textbf{student} or \textbf{local}.
A simplified conceptual matrix is:

\begin{center}
\begin{tabular}{>{\ttfamily}p{2.5cm} p{8.2cm}}
\toprule
Column & Meaning \\
\midrule
Identification & unique ID \\
My status & student or local text \\
My profile & own type / profile \\
My hobbies and interests & own hobbies \\
Hobbies and interests I/we appreciate & preferred hobbies in partner \\
My gender identity & own gender \\
Friend's gender identity & acceptable partner genders \\
The must match criteria 1/2 & hard criteria names \\
Specs for my criteria 1/2 & acceptable values \\
\bottomrule
\end{tabular}
\end{center}

Think of the input as a respondent table from which the program builds a \textbf{pairwise compatibility matrix}.
\end{frame}

\begin{frame}{How the scores are computed}
For each feasible pair, the program computes component similarities in $[0,1]$.
Typical examples:
\begin{itemize}
  \item \textbf{Profile:} Jaccard overlap of profile categories
  \item \textbf{Hobbies:} own hobbies vs partner's appreciated hobbies, both directions
  \item \textbf{Languages:} overlap of native and usable languages
  \item \textbf{Age:} soft penalty by distance, reduced if outside stated preference range
\end{itemize}
Final score:
\[
s_{ij} = \sum_k w_k c_{ij}^{(k)}
\]
where $w_k$ are normalized Chebychev weights and $c_{ij}^{(k)}$ are component similarities.
\end{frame}

\begin{frame}{How to read the results}
\textbf{1. Matches file:} \texttt{matches.csv}
\begin{itemize}
  \item one selected pair per row
  \item contains IDs and component scores
  \item \texttt{score} is the final weighted compatibility
\end{itemize}

\textbf{2. Diagnostics file:} \texttt{matches.diagnostics.csv}
\begin{itemize}
  \item one candidate pair per row
  \item tells which hard constraints passed or failed
  \item useful for explaining why feasibility is low
\end{itemize}

\textbf{3. Terminal statistics:}
\begin{itemize}
  \item number of students, locals, candidate pairs
  \item pass rates of each hard constraint
  \item unmatched counts after optimization
  \item min / avg / max score of selected pairs
\end{itemize}
\end{frame}

\begin{frame}[fragile]{Typical workflow}
\begin{enumerate}
  \item Prepare the Excel input table.
  \item Run the program:
\begin{lstlisting}[language=bash]
python main.py --excel "Refined Attributes with options.xlsx" \
  --objective lexi --verbose
\end{lstlisting}
  \item Inspect terminal statistics.
  \item Open \texttt{matches.csv} for final pairs.
  \item Open diagnostics if too many people remain unmatched.
  \item Adjust hard constraints or weights and rerun.
\end{enumerate}

\vspace{0.4em}
\textbf{Rule of thumb:}
\begin{itemize}
  \item too few feasible pairs $\Rightarrow$ soften hard constraints
  \item weak match quality $\Rightarrow$ adjust weights or objective priority
\end{itemize}
\end{frame}

\begin{frame}{Take-home message}
\begin{itemize}
  \item The solver separates the problem into:
  \begin{enumerate}
    \item \textbf{feasibility filtering} by hard constraints,
    \item \textbf{quality scoring} by weighted components,
    \item \textbf{ILP optimization} over feasible pairs.
  \end{enumerate}
  \item Hard constraints control \textbf{who may be matched}.
  \item Chebychev weights control \textbf{what a good match means}.
  \item \texttt{matches.csv} tells you the final solution.
  \item \texttt{diagnostics.csv} explains why some matches were impossible.
\end{itemize}
\end{frame}

\end{document}
